\section*{Введение}
\addcontentsline{toc}{section}{Введение}

Суперкомпьютерные системы используются для выполнения вычислительных задач, требующих значительного объёма процессорного времени и памяти. Такие системы обычно имеют кластерную архитектуру: множество вычислительных узлов объединены сетью и управляются специализированной системой диспетчеризации заданий. Одной из наиболее распространённых систем такого класса является SLURM.

Качество планирования в вычислительном кластере зависит от того, насколько точно пользователь оценивает требуемое время выполнения задания. На практике пользователь часто указывает время с запасом: фактическое время выполнения оказывается меньше запрошенного. Из-за этого диспетчер вынужден резервировать ресурсы на большее время, чем действительно необходимо, что усложняет планирование очереди и может снижать загрузку узлов.

Цель работы~--- развернуть учебную среду SLURM и разработать генератор пользовательских заданий, который воспроизводит свойства выбранного фрагмента реальной очереди. Для этого необходимо построить кластер из виртуальных машин, настроить сетевое взаимодействие между узлами, подготовить конфигурацию SLURM, проанализировать исходный набор данных и реализовать запуск сгенерированных заданий.

В отчёте описаны этапы подготовки виртуальных машин, настройки сетевой связности и SSH-доступа, установки MUNGE и SLURM, выбора параметров конфигурации, обработки исходного набора данных, генерации заданий и анализа результатов выполнения.
